% ===================================================================
%  CADRO N.º 4 — Idade madura e vellez.
%  Traduction 2026 d'apres texto/io, controlee sur texto/fr, texto/es
%  et texto/pt. Consigne : voir texto/gl/10-cadro-01.tex.
%
%  LA PARENTE PAR ALLIANCE EST LE TROISIEME LIEU OU LE GALICIEN
%  TRANCHE, apres le corps au tableau 2 et les jours au tableau 3.
%  L'espagnol et le portugais preposent tous deux un element au nom :
%  « suegro », « yerno », « cuñado » d'un cote, « sogro », « genro »,
%  « cunhado » de l'autre. Le galicien fait comme eux — « sogro »,
%  « xenro », « cuñado » — mais avec son X : « xenro » n'est ni
%  « yerno » ni « genro », et il ne peut se copier ni de l'un ni de
%  l'autre sans faute d'orthographe.
%
%  LES MALADIES DU SECOND TABLEAU, ELLES, SONT PAREILLES PARTOUT, et
%  c'est le contraire du corps au tableau 2. « Reumatismo »,
%  « pulmonía », « febre », « sincope », « catarro », « poción » : le
%  corps s'apprend a la maison, mais la maladie s'apprend chez le
%  medecin, et le medecin parle latin dans les trois langues. Le meme
%  partage qu'a la mer irlandaise, dans un tableau de famille.
% ===================================================================

%%K t04-apar-1 apar
\VUsaut{4.0mm}
\VUtitre{15.0pt}{60}{CADRO N.\textsuperscript{o} 4}
\VUsaut{1.6mm}
\VUtitre{11.4pt}{0}{Idade madura e vellez.}
\VUsaut{3.2mm}

%%K t04-tit-1 sub
\VUcentre{11.4pt}{0}{\textit{Primeira escena.}}
\VUsaut{1.6mm}
\VUtitre{11.4pt}{0}{O banquete de vodas.}
\VUsaut{2.6mm}

%%K t04-tit-2 sub
\VUpk{10.2pt}{40}{O outono.}
\VUsaut{3.0mm}

%%K t04-c1-01-1 p
1. --- Os mozos medraron. Un deles, Xoán, cásase. Asistimos ao \VUgras{banquete de
vodas}, que xunta arredor da mesma mesa as familias dos noivos.

%%K t04-c1-02-1 p
2. --- Estamos no \VUgras{outono,} no \VUgras{mes de setembro.} O vento frío de
\VUgras{outubro} e de \VUgras{novembro} aínda non espiu o campo da súa follaxe verde.
O aire do serán é morno e recendente; por iso a cea faise baixo o \VUgras{alpendre}
\textsuperscript{(8)}, ao pé do xardín.

%%K t04-c1-03-1 p
3. --- Lonxe, no \VUgras{outeiro} \textsuperscript{(3)}, está a casa dos pais de
Xoán, un elegante \VUgras{pazo} \textsuperscript{(1)} agochado na verdura; belas
\VUgras{viñas}, que dan un viño excelente, cobren a costa do outeiro. O viticultor
cultívaas e coídaas.

%%K t04-c1-04-1 p
4. --- Por riba das viñas pasa un bando de \VUgras{perdices} \textsuperscript{(2)}
asustadas pola chegada do \VUgras{gardabosques} \textsuperscript{(5)}. Este, coa
\VUgras{escopeta} baixo o brazo e o \VUgras{zurrón de caza en bandoleira}, percorre
a \VUgras{fraga} \textsuperscript{(4)} co seu \VUgras{can} \textsuperscript{(6)}.
Vixía a propiedade e impide que os furtivos destrúan a caza.

%%K t04-c1-05-1 p
5. --- Fóra do \VUgras{alpendre} trepa unha \VUgras{parreira} da que penduran
mestos \VUgras{acios de uvas} \textsuperscript{(7)}.

%%K t04-c1-06-1 p
6. --- As belas flores de outono que adornan a \VUgras{mesa} \textsuperscript{(16)}
son \VUgras{crisantemos} \textsuperscript{(9)}; o \VUgras{pai}
\textsuperscript{(10)} da noiva puxo un deles na botoeira do seu frac.

%%K t04-c1-07-1 p
7. --- A comida está a rematar. O \VUgras{criado} \textsuperscript{(11)} comeza a
traer os pratos da sobremesa e axiña vai retirar as distintas pezas do servizo,
\VUgras{culleres} \textsuperscript{(14)}, \VUgras{garfos} \textsuperscript{(12)},
\VUgras{coitelos} \textsuperscript{(13)}, \VUgras{fontes} \textsuperscript{(15)}, que
xa non fan falta.

%%K t04-tit-3 sub
\VUpk{10.2pt}{40}{A parentela.}
\VUsaut{3.0mm}

%%K t04-c1-08-1 p
8. --- Todos parecen ledos, e o sorriso co que a \VUgras{nai} \textsuperscript{(17)}
do noivo mira os seus fillos proba a súa felicidade.

%%K t04-c1-09-1 p
9. --- Esa voda une dúas familias ricas do país. Xoán, o \VUgras{esposo}
\textsuperscript{(18)}, é \VUgras{fillo} dun gran propietario e faise o
\VUgras{xenro} dun hábil industrial.

%%K t04-c1-10-1 p
10. --- Os \VUgras{esposos} son novos e, aínda así, hai moito que estaban
\VUgras{prometidos.} A \VUgras{muller nova} \textsuperscript{(19)}, Marta, está
encantadora co seu \VUgras{vestido branco} adornado con flores de laranxeira.

%%K t04-c1-11-1 p
11. --- O seu \VUgras{sogro} \textsuperscript{(20)} está moi feliz de ter unha
\VUgras{nora} tan xentil, e a \VUgras{sogra} \textsuperscript{(21)} de Xoán sorrí ao
ver tan bela a súa \VUgras{filla.}

%%K t04-c1-12-1 p
12. --- Na punta da mesa sentan os demais \VUgras{convidados} \textsuperscript{(22)},
\VUgras{parentes, amigos, curmáns, curmás.} Un \VUgras{mestresala}
\textsuperscript{(23)}, cun pano de servir baixo o brazo, sérveos con dilixencia.

%%K t04-tit-4 sub
\VUpk{10.2pt}{40}{Os amigos.}
\VUsaut{3.0mm}

%%K t04-c1-13-1 p
13. --- Á dereita da mesa, velaquí unha \VUgras{señorita} \textsuperscript{(24)},
\VUgras{amiga} da nova esposa, e despois o \VUgras{irmán} desta, que é o seu
\VUgras{padriño de honra} \textsuperscript{(25)}. Parola coa súa \VUgras{madriña de
honra} \textsuperscript{(26)}, irmá do esposo, que, con esa voda, faise a
\VUgras{cuñada} de Marta. É unha moza encantadora. Ten belas \VUgras{xoias, un
colar de perlas} e unha \VUgras{pulseira} de ouro.

%%K t04-c1-14-1 p
14. --- Preto deles, o señor que está de pé e ergue o seu vaso é un \VUgras{amigo}
\textsuperscript{(27)} da familia. É unha das \VUgras{testemuñas do esposo.}
\VUgras{Brinda,} bebe pola \VUgras{saúde} dos novos esposos e deséxalles unha longa
felicidade. Vai vestido con traxe de cerimonia, con frac e \VUgras{zapatos de
charol} \textsuperscript{(28)}. É o acompañante da \VUgras{avoa}
\textsuperscript{(29)}, que é \VUgras{viúva.}

%%K t04-c1-15-1 p
15. --- O mozo de \VUgras{frac} \textsuperscript{(31)}, de bigote negro e fino, é o
\VUgras{cuñado} \textsuperscript{(30)} de Xoán. É un enxeñeiro eminente. É o veciño
dunha \VUgras{dama nova} \textsuperscript{(33)} moi elegante, a \VUgras{irmá maior}
de Marta e a nai da meniña garrida que vén, dando o brazo ao seu curmán, ofrecer unha
flor e \VUgras{desexarlle} \VUgras{boas tardes.} Esa dama ten uns \VUgras{pendentes}
moi belos e unha elegante \VUgras{touca.}

%%K t04-c1-16-1 p
16. --- O curmán pequeno, tan raro coa súa \VUgras{blusa} \textsuperscript{(36)} e a
súa \VUgras{manga} \textsuperscript{(37)} inchada, dá moita mágoa. É \VUgras{orfo}
\textsuperscript{(35)}. Moi novo perdeu o pai e a nai. O seu tío é o seu
\VUgras{titor} \textsuperscript{(38)} e ficou solteiro para educar o pobre neno. É un
home bondadoso. Está algo calvo e é moi míope; usa \VUgras{lentes.}

%%K t04-tit-5 sub
\VUsaut{2.4mm}
\VUpk{10.2pt}{40}{A sobremesa.}
\VUsaut{3.0mm}

%%K t04-c1-17-1 p
17. --- Detrás dos convidados, nunha mesa cuberta cun \VUgras{mantel,} está preparada
a \VUgras{sobremesa} \textsuperscript{(39)}. Nunha copa grande, velaquí froitas,
\VUgras{pexegos} \textsuperscript{(40)}, \VUgras{peras} \textsuperscript{(41)},
\VUgras{mazás} \textsuperscript{(42)}, \VUgras{albaricoques} \textsuperscript{(43)} e
\VUgras{uvas} \textsuperscript{(44)}. Preto, \VUgras{ananases}
\textsuperscript{(45)}, un fermoso \VUgras{bolo} \textsuperscript{(46)},
\VUgras{botellas de licor} \textsuperscript{(48)} e \VUgras{champaña}
\textsuperscript{(49)}, e tamén columnas de \VUgras{pratos} \textsuperscript{(47)}
limpos.

%%K t04-c1-18-1 p
18. --- O \VUgras{adegueiro} \textsuperscript{(50)} descorcha unha \VUgras{botella}
\textsuperscript{(52)} de viño vello cun \VUgras{sacarrollas}
\textsuperscript{(51)}. A \VUgras{rolla de cortiza} \textsuperscript{(53)} parece moi
difícil de sacar. Ese adegueiro ten un \VUgras{pano de servir}
\textsuperscript{(54)} baixo o brazo.

%%K t04-c1-19-1 p
19. --- Despois da cea, os señores irán ao fumadoiro e as damas irán prepararse para o
baile.

%%K t04-tit-6 sub
\VUsaut{5.0mm}
\VUcentre{11.4pt}{0}{\textit{Segunda escena.}}
\VUsaut{1.6mm}
\VUpk{11.4pt}{40}{O santo do avó.}
\VUsaut{2.6mm}

%%K t04-tit-7 sub
\VUpk{10.2pt}{40}{A visita.}
\VUsaut{3.0mm}

%%K t04-c2-01-1 p
1. --- Pasaron dez anos; os pais de Xoán fixéronse vellos e queren moito os seus tres
netos, dous \VUgras{netos} \textsuperscript{(2)} e unha \VUgras{neta}
\textsuperscript{(3)}. Como hoxe é \VUgras{o santo do avó,} os fillos veñen desexarlles
un bo día de festa.

%%K t04-c2-02-1 p
2. --- O pequeno Pedro é aínda moi novo; ten só tres anos. Ve achegarse o
\VUgras{gato} \textsuperscript{(1)}. Quere acariñar o seu belo \VUgras{pelo} sedoso e
o seu longo \VUgras{rabo}; non pensa que baixo as graciosas \VUgras{patas} se agochan
unllas maliciosas e agudas que quizais o rabuñen.

%%K t04-c2-03-1 p
3. --- A súa irmá Gabriela, que lle dá a man, é moito máis seria, e Marcelo, coa súa
grande \VUgras{capa} \textsuperscript{(4)} e o seu \VUgras{cinto} de coiro
\textsuperscript{(5)}, semella algo feito home, malia a súa manga curta, que deixa ver
as súas \VUgras{medias} \textsuperscript{(6)}.

%%K t04-c2-04-1 p
4. --- O pai deles puxo o seu groso \VUgras{sobretodo} \textsuperscript{(7)} de
inverno con \VUgras{colo de pel} \textsuperscript{(8)} e, no seu \VUgras{peto}
\textsuperscript{(9)}, leva un pano. A súa muller ten o seu sombreiro de feltro
adornado con \VUgras{plumas} \textsuperscript{(10)}, a súa \VUgras{chaqueta}
\textsuperscript{(11)}, a súa \VUgras{boa de pel} \textsuperscript{(12)} e o seu
\VUgras{manguito} \textsuperscript{(13)}.

%%K t04-c2-05-1 p
5. --- Fai moito frío, pois reina o inverno. A \VUgras{neve} \textsuperscript{(19)}
caeu durante todo o día e os tellados das casas están brancos de todo. Ao mesmo tempo
ca a \VUgras{noite,} o frío faise aínda máis agudo. No ceo a \VUgras{lúa}
\textsuperscript{(17)} e as \VUgras{estrelas} \textsuperscript{(18)} brillan e
atravesan \VUgras{a escuridade.}

%%K t04-tit-8 sub
\VUsaut{4.2mm}
\VUpk{10.2pt}{40}{A doenza.}
\VUsaut{3.0mm}

%%K t04-c2-06-1 p
6. --- O pobre \VUgras{cego} \textsuperscript{(20)} que cruza a praza diante da
\VUgras{farmacia} \textsuperscript{(16)}, guiado polo seu \VUgras{can de augas}
\textsuperscript{(21)}, é ben desgraciado. Xustina, a \VUgras{criada}
\textsuperscript{(14)}, trae da cociña unha \VUgras{cunca grande}
\textsuperscript{(15)} de \VUgras{tisana} moi quente e azucrada con xenerosidade.

%%K t04-c2-07-1 p
7. --- A \VUgras{avoa} \textsuperscript{(23)}, que quere buscar as súas
\VUgras{lentes de man} \textsuperscript{(26)} na mesa, para ler a \VUgras{receita}
\textsuperscript{(27)} que o \VUgras{médico} \textsuperscript{(25)} acaba de
escribir, érguese do seu sillón apoiándose no seu \VUgras{bastón}
\textsuperscript{(24)}.

%%K t04-c2-08-1 p
8. --- A miúdo sofre de \VUgras{achaques, mareos, catarros, dores de dentes}; as
súas pernas están fracas e os seus ollos xa non son moi bos. Malia iso, de boa gana
búrlase do seu vello \VUgras{médico,} que sempre lle quere receitar unha
\VUgras{poción} \textsuperscript{(28)}. Hoxe está moi ledo porque ten arredor de si a
súa familia nova.

%%K t04-c2-09-1 p
9. --- Estase ben no cuarto ben pechado. Na \VUgras{lareira} \textsuperscript{(29)}
de mármore arde un \VUgras{lume fermosísimo} \textsuperscript{(30)}, e séntese un
feliz na doce \VUgras{luz} da \VUgras{lámpada} \textsuperscript{(31)} que acaban de
acender.

%%K t04-tit-9 sub
\VUsaut{4.2mm}
\VUpk{10.2pt}{40}{O avó.}
\VUsaut{3.0mm}

%%K t04-c2-10-1 p
10. --- Mesmo o \VUgras{avó} \textsuperscript{(33)} esqueceu que estaba doente, que o
seu \VUgras{reumatismo} o prende ao seu \VUgras{sillón} \textsuperscript{(34)} e que,
onte aínda, tivo \VUgras{febre e lipotimias.} Xa non lembra que o inverno pasado
sufriu de \VUgras{pulmonía} e que unha \VUgras{tose} rouca e penosa o fixo padecer
moito.

%%K t04-c2-10-2 p
E, aínda así, moi a duras penas se curou. Desde entón está algo mellor. Pero a perna
que apoia nun \VUgras{coxín} \textsuperscript{(38)} posto nun pequeno
\VUgras{tallo} \textsuperscript{(39)} tamén lle doe.

%%K t04-c2-11-1 p
11. --- Cando está coidadosamente envolto no seu quente \VUgras{roupón}
\textsuperscript{(35)} con grosos \VUgras{botóns} \textsuperscript{(36)} de nácar, e
cando ten preto de si, para lerlle o xornal, a pequena \VUgras{orfa}
\textsuperscript{(40)} vestida cunha \VUgras{cofia} branca, cun \VUgras{vestido}
azul \textsuperscript{(42)} e cunha \VUgras{esclavina} \textsuperscript{(41)}, non se
queixa.

%%K t04-c2-12-1 p
12. --- Cada ano, cando o \VUgras{calendario} \textsuperscript{(43)} marca de novo a
\VUgras{data} do primeiro de \VUgras{decembro,} agarda con impaciencia polos seus
\VUgras{netos,} pois sabe que non esquecerán esa \VUgras{data} importante.

%%K t04-c2-13-1 p
13. --- Lembra con emoción o tempo moi antigo en que el mesmo veu desexarlle un bo
santo ao glorioso \VUgras{mariscal} imperial, cuxo \VUgras{retrato}
\textsuperscript{(44)} pendura na parede e que é o \VUgras{bisavó} dos seus fillos.
